% GENERATED FILE. Edit canonical lessons, then run npm run generate:books.
% canonical-source-hash: fnv1a64:87f18aaee28fb4c1
% canonical-lessons: MR-C493-kubdi, MR-C493-cakkar-yene, MR-C493-adkhalne, MR-C493-kudkudne, MR-C493-langadne

\chapter{Dizzy and Stumbling}
\label{ch:mr-dizzy-and-stumbling}
\hlchaptermodality{493}

\begin{chapteropening}
\textbf{By the end of this chapter:} \emph{I can talk about dizziness, stumbling, shivering and limping.}

Complete the last lesson of chapter 493: I can talk about dizziness, stumbling, shivering and limping.
\end{chapteropening}

\section[\texorpdfstring{\mr{कुबडी} (kubḍī)}{कुबडी (kubḍī)}]{\mr{कुबडी} (kubḍī) — a crutch}

\label{lesson:MR-C493-kubdi}



\begin{quote}
Before the new one: say the Marathi for to get breathless, then the Marathi for to choke on food.
\end{quote}

\subsection*{You'll want to know: \mr{कुबडी}}
\textbf{\mr{कुबडी}} — \emph{kubḍī} — \textquotedblleft{}a crutch\textquotedblright{}.

\begin{quote}
\mr{पाय} \mr{मोडल्यामुळे} \mr{तो} \mr{कुबडी} \mr{घेऊन} \mr{चालतो}.
\end{quote}

\begin{quote}
\emph{pāy moḍlyāmuḷe to kubḍī gheūn cālto.}
\end{quote}

\begin{quote}
Since he broke his leg he walks with a crutch.
\end{quote}

Feminine. Plural: \mr{कुबड्या}, \emph{kubḍyā}.

Plural \mr{कुबड्या}, a pair of crutches.

\begin{cousinweb}
\textbf{\mr{कुबडी}} (\emph{kubḍī}) is a Marathi noun whose origin is not settled.
\end{cousinweb}

\begin{grammarlens}[title={What you've built}]
One more word for describing what you lived and felt.
\end{grammarlens}

\subsection*{Your turn}
\begin{itemize}
\raggedright
  \item \emph{Say it:} \emph{kubḍī}
  \item \emph{Say it:} \emph{kubḍī}, once more
  \item \emph{Say it:} \emph{pāy moḍlyāmuḷe to kubḍī gheūn cālto.} — Since he broke his leg he walks with a crutch.
  \item \emph{Recall:} say the Marathi for to get breathless, then the Marathi for to choke on food, then say \emph{kubḍī} again
\end{itemize}

\begin{tcolorbox}[breakable,colback=teal!4,colframe=teal!35!black,title={Before you move on}]
What does \mr{कुबडी} mean? (\textquotedblleft{}A crutch\textquotedblright{}.) Say it once more.
\end{tcolorbox}
\section[\texorpdfstring{\mr{चक्कर} \mr{येणे} (cakkar yeṇe)}{चक्कर येणे (cakkar yeṇe)}]{\mr{चक्कर} \mr{येणे} (cakkar yeṇe) — to feel dizzy}

\label{lesson:MR-C493-cakkar-yene}



\begin{quote}
Before the new one: say the Marathi for to choke on food, then the Marathi for a crutch.
\end{quote}

\subsection*{You'll want to know: \mr{चक्कर} \mr{येणे}}
\textbf{\mr{चक्कर} \mr{येणे}} — \emph{cakkar yeṇe} — \textquotedblleft{}to feel dizzy\textquotedblright{}.

\begin{quote}
\mr{उन्हात} \mr{उभं} \mr{राहिल्यामुळे} \mr{तिला} \mr{चक्कर} \mr{आली}.
\end{quote}

\begin{quote}
\emph{unhāt ubhaṁ rāhilyāmuḷe tilā cakkar ālī.}
\end{quote}

\begin{quote}
She felt dizzy from standing in the sun.
\end{quote}

Past (to him …): \mr{त्याला} \mr{चक्कर} \mr{आली}, \emph{tyālā cakkar ālī}.

Dative. \mr{चक्कर} also means a round: \mr{चक्कर} \mr{मारणे}, to go for a spin.

\begin{cousinweb}
\textbf{\mr{चक्कर} \mr{येणे}} (\emph{cakkar yeṇe}) is a Marathi phrase built on \mr{चक्कर}, a loan from Hindustani \emph{cakkar}, \textquotedblleft{}a whirl\textquotedblright{}, which goes back to Sanskrit \emph{cakra}, \textquotedblleft{}wheel\textquotedblright{} (CDIAL 4538); \mr{येणे} is the light verb.
\end{cousinweb}

\begin{grammarlens}[title={What you've built}]
One more word for describing what you lived and felt.
\end{grammarlens}

\subsection*{Your turn}
\begin{itemize}
\raggedright
  \item \emph{Say it:} \emph{cakkar yeṇe}
  \item \emph{Say it:} \emph{cakkar yeṇe}, once more
  \item \emph{Say it:} \emph{unhāt ubhaṁ rāhilyāmuḷe tilā cakkar ālī.} — She felt dizzy from standing in the sun.
  \item \emph{Recall:} say the Marathi for to choke on food, then the Marathi for a crutch, then say \emph{cakkar yeṇe} again
\end{itemize}

\begin{tcolorbox}[breakable,colback=teal!4,colframe=teal!35!black,title={Before you move on}]
What does \mr{चक्कर} \mr{येणे} mean? (\textquotedblleft{}To feel dizzy\textquotedblright{}.) Say it once more.
\end{tcolorbox}
\section[\texorpdfstring{\mr{अडखळणे} (aḍkhaḷṇe)}{अडखळणे (aḍkhaḷṇe)}]{\mr{अडखळणे} (aḍkhaḷṇe) — to stumble}

\label{lesson:MR-C493-adkhalne}



\begin{quote}
Before the new one: say the Marathi for a crutch, then the Marathi for to feel dizzy.
\end{quote}

\subsection*{You'll want to know: \mr{अडखळणे}}
\textbf{\mr{अडखळणे}} — \emph{aḍkhaḷṇe} — \textquotedblleft{}to stumble\textquotedblright{}.

\begin{quote}
\mr{दगडाला} \mr{पाय} \mr{लागून} \mr{तो} \mr{अडखळला}.
\end{quote}

\begin{quote}
\emph{dagḍālā pāy lāgūn to aḍkhaḷlā.}
\end{quote}

\begin{quote}
His foot hit a stone and he stumbled.
\end{quote}

Past (he …): \mr{तो} \mr{अडखळला}, \emph{to aḍkhaḷlā}.

Of feet and of words: \mr{बोलताना} \mr{अडखळणे}, to stumble over words.

\begin{cousinweb}
\textbf{\mr{अडखळणे}} (\emph{aḍkhaḷṇe}) is a Marathi verb whose origin is not settled.
\end{cousinweb}

\begin{grammarlens}[title={What you've built}]
One more word for describing what you lived and felt.
\end{grammarlens}

\subsection*{Your turn}
\begin{itemize}
\raggedright
  \item \emph{Say it:} \emph{aḍkhaḷṇe}
  \item \emph{Say it:} \emph{aḍkhaḷṇe}, once more
  \item \emph{Say it:} \emph{dagḍālā pāy lāgūn to aḍkhaḷlā.} — His foot hit a stone and he stumbled.
  \item \emph{Recall:} say the Marathi for a crutch, then the Marathi for to feel dizzy, then say \emph{aḍkhaḷṇe} again
\end{itemize}

\begin{tcolorbox}[breakable,colback=teal!4,colframe=teal!35!black,title={Before you move on}]
What does \mr{अडखळणे} mean? (\textquotedblleft{}To stumble\textquotedblright{}.) Say it once more.
\end{tcolorbox}
\section[\texorpdfstring{\mr{कुडकुडणे} (kuḍkuḍṇe)}{कुडकुडणे (kuḍkuḍṇe)}]{\mr{कुडकुडणे} (kuḍkuḍṇe) — to shiver with cold}

\label{lesson:MR-C493-kudkudne}



\begin{quote}
Before the new one: say the Marathi for to feel dizzy, then the Marathi for to stumble.
\end{quote}

\subsection*{You'll want to know: \mr{कुडकुडणे}}
\textbf{\mr{कुडकुडणे}} — \emph{kuḍkuḍṇe} — \textquotedblleft{}to shiver with cold\textquotedblright{}.

\begin{quote}
\mr{थंडीत} \mr{रस्त्यावरचा} \mr{कुत्रा} \mr{कुडकुडत} \mr{होता}.
\end{quote}

\begin{quote}
\emph{thaṇḍīt rastyāvarcā kutrā kuḍkuḍat hotā.}
\end{quote}

\begin{quote}
The street dog was shivering in the cold.
\end{quote}

Past (he …): \mr{तो} \mr{कुडकुडला}, \emph{to kuḍkuḍlā}.

Of cold, where \mr{थरथरणे} also covers fear.

\begin{cousinweb}
\textbf{\mr{कुडकुडणे}} (\emph{kuḍkuḍṇe}) is a Marathi verb built on a doubled syllable; its origin is not settled.
\end{cousinweb}

\begin{grammarlens}[title={What you've built}]
One more word for describing what you lived and felt.
\end{grammarlens}

\subsection*{Your turn}
\begin{itemize}
\raggedright
  \item \emph{Say it:} \emph{kuḍkuḍṇe}
  \item \emph{Say it:} \emph{kuḍkuḍṇe}, once more
  \item \emph{Say it:} \emph{thaṇḍīt rastyāvarcā kutrā kuḍkuḍat hotā.} — The street dog was shivering in the cold.
  \item \emph{Recall:} say the Marathi for to feel dizzy, then the Marathi for to stumble, then say \emph{kuḍkuḍṇe} again
\end{itemize}

\begin{tcolorbox}[breakable,colback=teal!4,colframe=teal!35!black,title={Before you move on}]
What does \mr{कुडकुडणे} mean? (\textquotedblleft{}To shiver with cold\textquotedblright{}.) Say it once more.
\end{tcolorbox}
\section[\texorpdfstring{\mr{लंगडणे} (laṅgaḍṇe)}{लंगडणे (laṅgaḍṇe)}]{\mr{लंगडणे} (laṅgaḍṇe) — to limp}

\label{lesson:MR-C493-langadne}



\begin{quote}
Before the new one: say the Marathi for to stumble, then the Marathi for to shiver with cold.
\end{quote}

\subsection*{You'll want to know: \mr{लंगडणे}}
\textbf{\mr{लंगडणे}} — \emph{laṅgaḍṇe} — \textquotedblleft{}to limp\textquotedblright{}.

\begin{quote}
\mr{पायाला} \mr{लागल्यामुळे} \mr{घोडा} \mr{लंगडत} \mr{होता}.
\end{quote}

\begin{quote}
\emph{pāyālā lāglyāmuḷe ghoḍā laṅgaḍat hotā.}
\end{quote}

\begin{quote}
The horse was limping from a hurt leg.
\end{quote}

Past (he …): \mr{तो} \mr{लंगडला}, \emph{to laṅgaḍlā}.

\mr{लंगडत} \mr{चालणे}, to walk with a limp.

\begin{cousinweb}
\textbf{\mr{लंगडणे}} (\emph{laṅgaḍṇe}) is a Marathi verb whose origin is not settled.
\end{cousinweb}

\begin{grammarlens}[title={What you've built}]
One more word for describing what you lived and felt.
\end{grammarlens}

\subsection*{Your turn}
\begin{itemize}
\raggedright
  \item \emph{Say it:} \emph{laṅgaḍṇe}
  \item \emph{Say it:} \emph{laṅgaḍṇe}, once more
  \item \emph{Say it:} \emph{pāyālā lāglyāmuḷe ghoḍā laṅgaḍat hotā.} — The horse was limping from a hurt leg.
  \item \emph{Recall:} say the Marathi for to stumble, then the Marathi for to shiver with cold, then say \emph{laṅgaḍṇe} again
\end{itemize}

\begin{tcolorbox}[breakable,colback=teal!4,colframe=teal!35!black,title={Before you move on}]
What does \mr{लंगडणे} mean? (\textquotedblleft{}To limp\textquotedblright{}.) Say it once more.
\end{tcolorbox}
